% =========================================================
% Topologies and Borel Generation
% =========================================================

\subsection{Topologies and Borel Generation}

\begin{exposition}
Topology is another closure structure on a family of subsets.  It differs from
sigma-algebras in exactly the way the later subjects need: topologies are
closed under arbitrary unions and finite intersections, while sigma-algebras
are closed under complements and countable unions.
\end{exposition}

\begin{definitionbox}{Definition (Topology)}
\begin{definition}[Topology]\label{def:topology}
Let \(X\) be a set.  A \emph{topology} on \(X\) is a family
\(\tau\subseteq\mathcal P(X)\) such that:
\begin{enumerate}[label=(\roman*)]
\item \(\varnothing\in\tau\);
\item \(X\in\tau\);
\item \(\tau\) is closed under arbitrary unions;
\item \(\tau\) is closed under finite intersections.
\end{enumerate}
The members of \(\tau\) are called \emph{open sets}.
\end{definition}
\end{definitionbox}

\begin{remark*}[Standard quantified statement]
\[
\tau\text{ is a topology on }X
\iff
\varnothing,X\in\tau,\quad
\tau\text{ is closed under arbitrary unions and finite intersections}.
\]
\end{remark*}

\begin{remark*}[Interpretation]
A topology is the set-family structure that remembers which subsets are open.
The arbitrary-union axiom lets open sets be assembled from local information;
the finite-intersection axiom lets finitely many local conditions be imposed at
the same time.
\end{remark*}

\begin{remark*}[Exposition]
A topology need not be closed under complements or countable intersections.
A sigma-algebra need not be closed under arbitrary unions.  Borel
sigma-algebras are the bridge: start from the open sets, then generate the
smallest sigma-algebra containing them.
\end{remark*}

\LeanFormalizes{def:topology}{lra-lean}{LRA.VolumeI.Set.Families}{TopologyOn}{pending}

\begin{dependencies}
\begin{itemize}
  \item \hyperref[def:closed-under-set-operation-vocabulary]{Closed Under Set Operations}
  \item \hyperref[def:indexed-union]{Indexed Union}
  \item \hyperref[def:indexed-intersection]{Indexed Intersection}
  \item \hyperref[thm:pairwise-closure-implies-finite-closure]{Pairwise Closure Implies Finite Closure}
\end{itemize}
\end{dependencies}

\begin{definitionbox}{Definition (Borel Sigma-Algebra)}
\begin{definition}[Borel Sigma-Algebra]\label{def:borel-sigma-algebra}
Let \((X,\tau)\) be a topological space.  The \emph{Borel sigma-algebra} of
\((X,\tau)\) is the sigma-algebra generated by the open sets:
\[
  \mathcal B(X,\tau)=\sigma(\tau).
\]
Its members are called \emph{Borel sets}.
\end{definition}
\end{definitionbox}

\begin{remark*}[Standard quantified statement]
\[
\mathcal B(X,\tau)=\sigma(\tau).
\]
\end{remark*}

\begin{remark*}[Interpretation]
The Borel sigma-algebra is the smallest sigma-algebra that contains every open
set.  It is the point where topology supplies the starting family and measure
theory supplies the closure rules.
\end{remark*}

\begin{remark*}[Exposition]
This chapter uses Borel generation tactically.  It does not pursue the Borel
hierarchy, cardinality of Borel sets, analytic sets, or descriptive set theory.
\end{remark*}

\LeanFormalizes{def:borel-sigma-algebra}{lra-lean}{LRA.VolumeI.Set.Families}{BorelSigmaAlgebra}{pending}

\begin{dependencies}
\begin{itemize}
  \item \hyperref[def:topology]{Topology}
  \item \hyperref[def:generated-sigma-algebra]{Generated Sigma-Algebra}
  \item \hyperref[def:sigma-algebra]{Sigma-Algebra}
\end{itemize}
\end{dependencies}

\begin{theorembox}{Theorem (Borel Sigma-Algebra Closure Laws)}
\begin{theorem}[Borel Sigma-Algebra Closure Laws]\label{thm:borel-sigma-algebra-closure-laws}
Let \((X,\tau)\) be a topological space.  Then:
\begin{enumerate}[label=(\roman*)]
\item every open set is Borel, equivalently \(\tau\subseteq\mathcal B(X,\tau)\);
\item \(\mathcal B(X,\tau)\) is a sigma-algebra on \(X\);
\item if \(\mathcal M\) is a sigma-algebra on \(X\) and
      \(\tau\subseteq\mathcal M\), then
      \(\mathcal B(X,\tau)\subseteq\mathcal M\).
\end{enumerate}
\hyperref[prf:borel-sigma-algebra-closure-laws]{Go to proof.}
\end{theorem}
\end{theorembox}

\begin{remark*}[Standard quantified statement]
\[
\tau\subseteq\mathcal B(X,\tau),\quad
\mathcal B(X,\tau)\text{ is a sigma-algebra on }X,\quad
\bigl(\mathcal M\text{ is a sigma-algebra on }X\ \wedge\ \tau\subseteq\mathcal M\bigr)
\to
\mathcal B(X,\tau)\subseteq\mathcal M.
\]
\end{remark*}

\begin{remark*}[Interpretation]
The Borel sigma-algebra is not merely a sigma-algebra containing the open sets;
it is the smallest such sigma-algebra.  This is the fact used whenever a proof
shows a property holds for all Borel sets by first checking that the property
defines a sigma-algebra containing the opens.
\end{remark*}

\LeanFormalizes{thm:borel-sigma-algebra-closure-laws}{lra-lean}{LRA.VolumeI.Set.Families}{TopologySubsetBorelSigmaAlgebra}{pending}
\LeanFormalizes{thm:borel-sigma-algebra-closure-laws}{lra-lean}{LRA.VolumeI.Set.Families}{BorelSigmaAlgebraIsSigmaAlgebra}{pending}
\LeanFormalizes{thm:borel-sigma-algebra-closure-laws}{lra-lean}{LRA.VolumeI.Set.Families}{BorelSigmaAlgebraMinimal}{pending}

\begin{dependencies}
\begin{itemize}
  \item \hyperref[def:borel-sigma-algebra]{Borel Sigma-Algebra}
  \item \hyperref[thm:generated-sigma-algebra-closure-laws]{Generated Sigma-Algebra Closure Laws}
  \item \hyperref[def:topology]{Topology}
  \item \hyperref[def:sigma-algebra]{Sigma-Algebra}
\end{itemize}
\end{dependencies}

\begin{definitionbox}{Definition (Metric Topology)}
\begin{definition}[Metric Topology]\label{def:metric-topology}
Let \(d:X\times X\to\mathbb R\) be a metric on \(X\).  For \(x\in X\) and
\(r>0\), the \emph{open ball} centered at \(x\) with radius \(r\) is
\[
  B_d(x,r)=\{y\in X:d(x,y)<r\}.
\]
A set \(U\subseteq X\) is \emph{\(d\)-open} if for every \(x\in U\) there is
some \(r>0\) such that
\[
  B_d(x,r)\subseteq U.
\]
The \emph{metric topology induced by \(d\)} is the family
\[
  \tau_d=\{U\subseteq X:U\text{ is }d\text{-open}\}.
\]
\end{definition}
\end{definitionbox}

\begin{remark*}[Standard quantified statement]
\[
\tau_d=\{U\subseteq X:\forall x\in U,\ \exists r>0,\ B_d(x,r)\subseteq U\}.
\]
\end{remark*}

\begin{remark*}[Predicate reading]
\[
  \tau_d\text{ is the metric topology induced by }d
  \iff
  \forall U\subseteq X,\ 
  U\in\tau_d\leftrightarrow
  \forall x\in U,\ \exists r>0,\ B_d(x,r)\subseteq U.
\]
\end{remark*}

\begin{remark*}[Interpretation]
A metric topology turns distance information into open-set information.  Once
the open sets are collected into \(\tau_d\), Borel generation proceeds exactly
as before:
\[
  \mathcal B(X,d)=\sigma(\tau_d).
\]
\end{remark*}

\LeanFormalizes{def:metric-topology}{lra-lean}{LRA.VolumeI.Set.Families}{DistanceFunction}{pending}
\LeanFormalizes{def:metric-topology}{lra-lean}{LRA.VolumeI.Set.Families}{MetricOn}{pending}
\LeanFormalizes{def:metric-topology}{lra-lean}{LRA.VolumeI.Set.Families}{MetricOpenBall}{pending}
\LeanFormalizes{def:metric-topology}{lra-lean}{LRA.VolumeI.Set.Families}{MetricOpenSet}{pending}
\LeanFormalizes{def:metric-topology}{lra-lean}{LRA.VolumeI.Set.Families}{MetricTopology}{pending}

\begin{dependencies}
\begin{itemize}
  \item \hyperref[def:topology]{Topology}
  \item \hyperref[def:borel-sigma-algebra]{Borel Sigma-Algebra}
  \item \hyperref[def:subset]{Subset}
\end{itemize}
\end{dependencies}

\begin{definitionbox}{Definition (Borel Sigma-Algebra on \(\mathbb R^n\))}
\begin{definition}[Borel Sigma-Algebra on \(\mathbb R^n\)]\label{def:borel-sigma-algebra-rn}
Let \(d_{\mathrm{Euc}}\) be the Euclidean metric on \(\mathbb R^n\), and let
\(\tau_{\mathrm{Euc}}\) be the metric topology induced by \(d_{\mathrm{Euc}}\).
The \emph{Borel sigma-algebra on \(\mathbb R^n\)} is
\[
  \mathcal B(\mathbb R^n)
  =
  \sigma(\tau_{\mathrm{Euc}}).
\]
Equivalently, \(\mathcal B(\mathbb R^n)\) is the smallest sigma-algebra on
\(\mathbb R^n\) containing every Euclidean open set.
\end{definition}
\end{definitionbox}

\begin{remark*}[Standard quantified statement]
\[
\mathcal B(\mathbb R^n)=\sigma(\tau_{\mathrm{Euc}}).
\]
\end{remark*}

\begin{remark*}[Interpretation]
This is the Borel structure used in multivariable analysis and measure theory.
The metric gives open balls, open balls determine Euclidean open sets, and the
Euclidean open sets generate the Borel sigma-algebra.
\end{remark*}

\begin{remark*}[Exposition]
The Euclidean metric formula and its metric-space properties belong to the
metric-space development.  This set-theory chapter only needs the set-family
construction:
\[
  d_{\mathrm{Euc}}
  \leadsto
  \tau_{\mathrm{Euc}}
  \leadsto
  \sigma(\tau_{\mathrm{Euc}}).
\]
\end{remark*}

\LeanFormalizes{def:borel-sigma-algebra-rn}{lra-lean}{LRA.VolumeI.Set.Families}{EuclideanSpace}{pending}
\LeanFormalizes{def:borel-sigma-algebra-rn}{lra-lean}{LRA.VolumeI.Set.Families}{BorelSigmaAlgebraRn}{pending}

\begin{dependencies}
\begin{itemize}
  \item \hyperref[def:metric-topology]{Metric Topology}
  \item \hyperref[def:borel-sigma-algebra]{Borel Sigma-Algebra}
  \item \hyperref[thm:borel-sigma-algebra-closure-laws]{Borel Sigma-Algebra Closure Laws}
\end{itemize}
\end{dependencies}
